\clearpage\thispagestyle{empty}
\pdfbookmark[0]{Quellen zum Juno-Bericht vom Februar 1805}{juno180502-sources}
\begin{center}\textbf{Quellen zum Juno-Bericht vom Februar 1805}\end{center}
\noindent Diese KI-Übertragung in heutiges Deutsch folgt dem Auszug aus Franz Xaver von Zachs Bericht in \href{https://gdz.sub.uni-goettingen.de/id/PPN236007467}{Gauß’ \textit{Werken}, Band~VI} (Göttingen, 1874), S.~260--261. Die entsprechende Passage in \textit{Ueber die Juno}, Beitrag XVII der \textit{Monatlichen Correspondenz}, Band~XI, Februar 1805, reicht von \href{https://www.digitale-sammlungen.de/en/view/bsb10538603?page=195}{S.~185} bis \href{https://www.digitale-sammlungen.de/en/view/bsb10538603?page=198}{S.~188}.
\par\smallskip\noindent Carl Friedrich Gauß berechnete die vierten Elemente der Juno aus eigenen Beobachtungen und Beobachtungen von Wilhelm Olbers. Der im Bericht besprochene Stern steht auf \href{https://archive.org/download/histoireclestef00lalagoog/page/n162.jpg}{S.~135} in Jérôme de Lalandes \textit{Histoire céleste française} (Paris, Jahr IX [1801]).
\par\smallskip\noindent Die Anmerkung der Werkausgabe verweist auf die \hyperlink{native-111}{auf S.~106 abgedruckten Vorschriften} und auf Gauß’ \hyperlink{native-124}{\textit{Der Zodiakus der Juno}}, erschienen im Märzheft auf \href{https://www.digitale-sammlungen.de/en/view/bsb10538603?page=237}{S.~225--228}.
\par\smallskip\noindent Übertragung und Satz wurden mit KI erstellt. Die Digitalisate der Zeitschrift stammen von der Bayerischen Staatsbibliothek.
\clearpage
